\documentclass[12pt]{article}

\usepackage[utf8]{inputenc}
\usepackage[T1]{fontenc}
\usepackage[polish]{babel}


\title{Dokumentacja funkcjonalna}
\author{Mateusz Kamieniecki \\ Jan Rękas}
\date{\today}


\begin{document}

\maketitle

\section{Cel projektu}
Celem projektu jest stworzenie aplikacji w języku C, umożliwiającej wyznaczanie
współrzędnych wierzchołków grafów planarnych w sposób czytelny i intuicyjny dla
użytkownika. Program pozwala na eksperymentowanie z dwoma różnymi algorytmami
rozmieszczania wierzchołków oraz generowanie wyników w formacie binarnym,
przeznaczonym do dalszego przetwarzania, lub w formacie tekstowym, przyjaznym
dla użytkownika.

\section{Opis funkcjonalności}

Program udostępnia następujące funkcje:

\begin{itemize}
	\item wczytywanie grafu z pliku tekstowego,
	\item wybór algorytmu wyznaczania współrzędnych wierzchołków,
	\item obliczanie współrzędnych wierzchołków zgodnie z wybranym algorytmem,
	\item zapis wyników w formacie binarnym lub tekstowym,
	\item wyświetlanie komunikatów o błędach w przypadku nieprawidłowego formatu danych wejściowych,
\end{itemize}

\section{Format danych wejściowych}

Graf wejściowy powinien być zapisany w pliku tekstowym, gdzie każda linia
reprezentuje krawędź grafu. Format linii powinien być następujący:

\begin{verbatim}
<nazwa_krawędzi> <wierzchołek_A> <wierzchołek_B> <waga_krawędzi>
...
\end{verbatim}

Przykład:
\begin{verbatim}
AB 1 2 1
BC 2 3 1
CD 3 4 1
DB 4 2 1.407
\end{verbatim}

\section{Format danych wyjściowych}

Program generuje plik zawierający współrzędne wierzchołków w formacie binarnym
lub tekstowym.

\subsection{Format tekstowy}
Plik wyjściowy w formacie tekstowym ma rozszerzenie .txt
Dane w formacie tekstowym są zapisywane w następujący sposób:

\begin{verbatim}
<liczba_wierzchołków>
<wierzchołek> <współrzędna_x> <współrzędna_y>
...
\end{verbatim}

Przykład:

\begin{verbatim}
4
1 0.0 0.0
2 1.0 0.0
3 1.0 1.0
4 0.0 1.0
\end{verbatim}

\subsection{Format binarny}

Dane w formacie binarnym są zapisywane jako ciąg bajtów. Na początku pliku
zapisana jest liczba wierzchołków (\texttt{int}, 4 bajty), a następnie każdy
wierzchołek jest reprezentowany przez trzy wartości: identyfikator wierzchołka
(4 bajty), współrzędna x (8 bajtów) i współrzędna y (8 bajtów). W sumie każdy
wierzchołek zajmuje 20 bajtów. Plik wyjściowy ma w formacie binarnym rozszerzenie .bin

Format danych binarnych można przedstawić następująco:

\begin{verbatim}
[liczba_wierzchołków]
[ID_wierzchołka][współrzędna_x][współrzędna_y]
...
\end{verbatim}

Przykład dla wierzchołków 1, 2, 3 i 4 z współrzędnymi (0.0, 0.0), (1.0, 0.0), (1.0, 1.0) i (0.0, 1.0):
\begin{verbatim}
[ 04 00 00 00 ]
[ 01 00 00 00 ][00 00 00 00 00 00 00 00][00 00 00 00 00 00 00 00]
[ 02 00 00 00 ][00 00 00 00 00 00 F0 3F][00 00 00 00 00 00 00 00]
[ 03 00 00 00 ][00 00 00 00 00 00 F0 3F][00 00 00 00 00 00 F0 3F]
[ 04 00 00 00 ][00 00 00 00 00 00 00 00][00 00 00 00 00 00 F0 3F]
\end{verbatim}
{\small\begin{verbatim}
[ 00 00 00 00 ] - liczba int zapisana w formacie little-endian
[00 00 00 00 00 00 00 00] - liczba double zapisana w formacie ISO 754
[00 00 00 00 00 00 00 00] - liczba double zapisana w formacie ISO 754
\end{verbatim}}

W tym przykładzie, identyfikatory wierzchołków są zapisane jako 4-bajtowe
liczby całkowite, a współrzędne x i y są zapisane jako 8-bajtowe liczby
zmiennoprzecinkowe (double). W rzeczywistości, dane binarne będą zapisane bez
spacji i nawiasów, a powyższy przykład jest jedynie ilustracją formatu danych.

\section{Flagi sterujące programem}

Program jest wywoływany z linii poleceń w następujący sposób:

\begin{verbatim}
./bin/planar_graph [flagi]
\end{verbatim}
Możliwe argumenty wywołania programu są opisane w poniższej tabelce to:

\begin{center}
	\begin{tabular}{|c|c|p{7cm}|c|}
		\hline
		\textbf{Flaga} & \textbf{Argument}          & \textbf{Opis}                                    & \textbf{Domyślnie}        \\
		\hline
		\texttt{-i}    & \texttt{<nazwa>}           & Nazwa pliku wejściowego z grafem                 & \texttt{"input.txt"}      \\
		\hline
		\texttt{-o}    & \texttt{<nazwa>}           & Ścieżka do pliku wyjściowego z wynikami          & \texttt{Katalog porgramu} \\
		\hline
		\texttt{-n}    & \texttt{<nazwa>}           & Nazwa pliku wyjściowego z wynikami               & \texttt{"output"}         \\
		\hline
		\texttt{-a}    & \texttt{<numer algorytmu>} & Wybór algorytmu rozmieszczania wierzchołków      & \texttt{1}                \\
		\hline
		\texttt{-b}    & —                          & Zapis wyników w formacie binarnym                & wyłączone*                \\
		\hline
		\texttt{-h}    & —                          & Zapis wyników w formacie czytelnym dla człowieka & wyłączone*                \\
		\hline
	\end{tabular}
\end{center}
* Jeżeli nie podano ani \texttt{-b}, ani \texttt{-h}, program automatycznie wybiera zapis binarny.

\subsection{Opis flag}

\begin{description}
	\item[\texttt{-i <ścieżka>}]
	      Określa ścieżkę do pliku wejściowego zawierającego opis grafu w formacie
	      tekstowym (opisanym w sekcji 3). Jeśli flaga nie zostanie podana, program
	      szuka pliku o nazwie \texttt{input.txt} w bieżącym katalogu roboczym.

	\item[\texttt{-o <ścieżka>}]
	      Określa ścieżkę do katalogu wyjściowego, w którym zostanie utworzony plik wyjściowy.
	      Jeśli flaga nie zostanie podana, wyniki zostaną zapisane w bieżącym katalogu.
	\item[\texttt{-n <nazwa>}]
	      Określa nazwę pliku wyjściowego. Jeśli flaga nie zostanie podana, wyniki zostaną zapisane w pliku o nazwie \texttt{output}.

	\item[\texttt{-a <nr>}]
	      Wybór algorytmu wyznaczania współrzędnych. Dostępne wartości:
	      \begin{itemize}
		      \item \texttt{1} — Tutte embedding
		      \item \texttt{2} — Spectral layout
	      \end{itemize}
	      Podanie nieobsługiwanej wartości skutkuje błędem i zakończeniem programu
	      (patrz sekcja wyjątki).

	\item[\texttt{-b}]
	      Włącza zapis wyników w formacie binarnym (.bin)(opisanym w sekcji 4.2).
	      Flaga nie przyjmuje dodatkowego argumentu.

	\item[\texttt{-h}]
	      Włącza zapis wyników w formacie tekstowym (.txt), czytelnym dla człowieka
	      (opisanym w sekcji 4.1). Flaga nie przyjmuje dodatkowego argumentu.
\end{description}

\subsection{Przykładowe wywołania}
Wczytanie grafu z pliku "graf.txt", zapis wyniku tekstowego do katalogu "wyniki", algorytm \texttt{alg1}
\begin{verbatim}
./bin/planar_graph -i graf.txt -o wyniki -n wynik -a alg1 -h
\end{verbatim}
Wczytanie grafu z pliku domyślnego, zapis binarny
\begin{verbatim}
./bin/planar_graph -b
\end{verbatim}
Zapis obu formatów jednocześnie
\begin{verbatim}
./bin/planar_graph -i graf.txt -n wynik -b -h
\end{verbatim}

\section{Sytuacje wyjątkowe}
\label{sec:wyjatki}

Poniżej opisano sytuacje wyjątkowe obsługiwane przez program. Komunikaty mogą być
wypisywane na \texttt{stderr} lub \texttt{stdout} (zgodnie z implementacją danej
funkcji), a kod powrotu zależy od przypadku i jest podany osobno dla każdego punktu.

\begin{description}

	\item[Nieznana lub niekompletna flaga]
	      \hfill \\
	      \textit{Sytuacja:} Podano flagę, której program nie rozpoznaje, lub flagę
	      wymagającą argumentu (np. \texttt{-i}, \texttt{-o}, \texttt{-n}, \texttt{-a}) bez
	      podania wartości. \\
	      \textit{Komunikat:}
	      \begin{verbatim}
Nieznany lub niekompletny argument: <flaga>
	      \end{verbatim}
	      \textit{Kod powrotu:} \texttt{1}

	\item[Nieobsługiwany algorytm]
	      \hfill \\
	      \textit{Sytuacja:} Wartość podana po fladze \texttt{-a} nie odpowiada
	      żadnemu z dostępnych algorytmów. \\
	      \textit{Komunikat:}
	      \begin{verbatim}
Taki algorytm nie jest obsługiwany
	      \end{verbatim}
	      \textit{Kod powrotu:} \texttt{1}

	\item[Nie można otworzyć pliku wejściowego]
	      \hfill \\
	      \textit{Sytuacja:} Plik wskazany flagą \texttt{-i} (lub domyślny plik
	      \texttt{input.txt}) nie istnieje lub brak uprawnień do odczytu. \\
	      \textit{Komunikat:}
	      \begin{verbatim}
Blad: nie mozna otworzyc pliku wejściowego: <nazwa_pliku>
	      \end{verbatim}
	      \textit{Kod powrotu:} \texttt{1}

	\item[Nie można otworzyć pliku wyjściowego]
	      \hfill \\
	      \textit{Sytuacja:} Program nie może utworzyć ani zapisać pliku wskazanego
	      flagą \texttt{-o} (np. brak uprawnień do zapisu w danym katalogu). \\
	      \textit{Komunikat:}
	      \begin{verbatim}
Nie udało się otworzyć pliku do zapisu: <systemowy_opis_błędu>
	      \end{verbatim}
	      \textit{Kod powrotu:} \texttt{0} (program wypisuje błąd zapisu, ale nie wymusza \texttt{EXIT\_FAILURE})

	\item[Nieprawidłowe dane wejściowe (brak dodatnich ID)]
	      \hfill \\
	      \textit{Sytuacja:} Plik wejściowy jest pusty albo nie zawiera żadnych dodatnich
	      identyfikatorów wierzchołków. \\
	      \textit{Komunikat:}
	      \begin{verbatim}
Blad: plik ma zły format, mozliwe: 1 lub 0 wierzcholkow, same niedodatnie id wierzcholkow.
	      \end{verbatim}
	      \textit{Kod powrotu:} \texttt{1}

	\item[Ostrzeżenie: krawędź z niedodatnim ID]
	      \hfill \\
	      \textit{Sytuacja:} W pliku wejściowym pojawia się krawędź z wierzchołkiem o ID $\le 0$. \\
	      \textit{Komunikat:}
	      \begin{verbatim}
Ostrzeżenie: Ignorowanie krawędzi z niedodatnim ID wierzchołka.
	      \end{verbatim}
	      \textit{Kod powrotu:} brak wpływu (ostrzeżenie, przetwarzanie trwa dalej)

	\item[Ostrzeżenie: krawędź poza zakresem ID]
	      \hfill \\
	      \textit{Sytuacja:} Dodawana krawędź odwołuje się do ID spoza zakresu \texttt{1..N}. \\
	      \textit{Komunikat:}
	      \begin{verbatim}
Ignorowanie krawędzi z nieprawidłowym ID wierzchołka: <from> -> <to> (dozwolony zakres: 1..<N>)
	      \end{verbatim}
	      \textit{Kod powrotu:} brak wpływu (ostrzeżenie, przetwarzanie trwa dalej)

	\item[Błąd alokacji pamięci]
	      \hfill \\
	      \textit{Sytuacja:} Program nie może zaalokować wymaganej pamięci podczas
	      wczytywania grafu lub wykonywania obliczeń. \\
	      \textit{Komunikat:}
	      \begin{verbatim}
Nie udało się zaalokować pamięci dla danych wektora
Nie udało się realokować pamięci dla danych wektora
Nie udało się zaalokować pamięci dla grafu
Nie udało się zaalokować pamięci dla list sąsiedztwa grafu
Nie udało się zaalokować pamięci dla nazwy krawędzi
Nie udało się zaalokować pamięci dla nazwy krawędzi odwrotnej
	      \end{verbatim}
	      \textit{Kod powrotu:} \texttt{1}

	\item[Problem numeryczny (algorytm Tutte)]
	      \hfill \\
	      \textit{Sytuacja:} W metodzie Gaussa pojawia się zerowy pivot. \\
	      \textit{Komunikat:}
	      \begin{verbatim}
Problem numeryczny: zerowy element główny
	      \end{verbatim}
	      \textit{Kod powrotu:} \texttt{0} (komunikat informacyjny)

	\item[Brak cyklu w grafie (algorytm Tutte)]
	      \hfill \\
	      \textit{Sytuacja:} Algorytm nie znajduje cyklu w grafie. \\
	      \textit{Komunikat:}
	      \begin{verbatim}
Nie znaleziono cyklu
	      \end{verbatim}
	      \textit{Kod powrotu:} \texttt{0}

\end{description}




\end{document}
